\section{Introduction}
\label{sec:introduction}

Large language models (LLMs) are sensitive to
how prompts are formulated, and even small
changes in wording can affect their predictions.
Errica et al.~\cite{Errica:2025}, in
\textit{What Did I Do Wrong? Quantifying LLMs'
Sensitivity and Consistency to Prompt Engineering},
investigate how different prompting strategies
and prompt rephrasings affect model sensitivity,
consistency, and classification performance
across multiple datasets.

\par\medskip

In this reproducibility study, we evaluate
whether the findings of the original paper
can be reproduced using Mixtral-8x7B and
Llama‐3‐70B‐Instruct. We additionally extend the
evaluation to Qwen3.5-27B and DeepSeek-R1 7B.
Following the original evaluation framework,
we compare three prompting strategies
(Simple, Detail, and 1-shot) across five
classification datasets (TREC, CB, RTE,
DBPedia, and WoS), measuring sensitivity,
consistency, and micro-F1.


\section{Scope of reproducibility}
\label{sec:claims}

Errica et al.~\cite{Errica:2025} investigate
how prompt engineering affects the sensitivity,
consistency, and classification performance
of large language models. Their experiments
compare different prompting strategies and
rephrasings across multiple classification
datasets.

Our reproducibility study focuses on the
following three testable claims:

\begin{itemize}
    \item \textbf{Claim 1: Prompt sensitivity.}
    LLM predictions are sensitive to prompt
    rephrasings, and the degree of sensitivity
    varies across prompting strategies
    and datasets.

    \item \textbf{Claim 2: Prediction consistency.}
    The consistency of LLM predictions across
    prompt rephrasings varies depending on
    the prompting strategy and dataset.

    \item \textbf{Claim 3: Classification performance.}
    Prompting strategies affect classification
    performance, measured by micro-F1, and
    no single strategy consistently achieves
    the highest performance across all datasets.
\end{itemize}

We test these claims by reproducing the
experiments for Mixtral-8x7B and Llama‐3‐70B‐Instruct
and extending the evaluation to Qwen3.5-27B
and DeepSeek-R1 7B.


\section{Methodology}
We reproduce the evaluation framework of the
original study using the authors' publicly
available code~\cite{original_code}.
Our primary objective is to reproduce the
reported results for Mixtral-8x7B-Instruct-v0.1,
followed by an evaluation of Llama‐3‐70B‐Instruct.
We additionally investigate Qwen3.5-27B
and DeepSeek-R1 7B to extend the analysis
to other language models.

The original evaluation procedures are retained
where possible, with modifications described
below to accommodate the available hardware,
execution constraints, and model-specific
requirements.\footnote{The original code was
developed by NEC Laboratories Europe GmbH.
Its license requires this acknowledgement.}

\paragraph{Mixtral-8x7B-Instruct-v0.1.}
\mbox{}\par
We reproduce the original evaluation using
Mixtral-8x7B-Instruct-v0.1
\cite{mixtral-8x7b-instruct-hf}, following
the authors' publicly available implementation.
The model is served through vLLM on the
ITU HPC DGX1 node. Minor adaptations were
introduced to accommodate computational
constraints and longer input sequences,
as detailed in the Mixtral experimental
setup section.

\paragraph{Llama‐3‐70B‐Instruct}
\mbox{}\par
[same as Mixtral]

\paragraph{Qwen3.5-27B.}
\mbox{}\par
To extend the original evaluation, we
additionally evaluate Qwen3.5-27B using
a local inference setup. The model is
registered in the authors' model table.
For CB, the main Qwen run uses the
training file that the CB notebook reads,
the same samples as our Mixtral run. A
second Qwen run uses the official test
samples with our rebuilt labels
(Section~\ref{sec:datasets}). The
truncated upstream TREC notebook is
restored.

\paragraph{DeepSeek-R1 7B.}
\mbox{}\par
We additionally evaluate DeepSeek-R1 7B,
a smaller reasoning model, on the ASUS1
node of the ITU HPC cluster. The model
is served through Ollama and registered
in the authors' model table. A wrapper
separates its reasoning text from the
final answer before applying the authors'
answer extraction procedure (Section~\ref{sec:setup}).

Computational resources and execution times
are reported in Section~\ref{sec:compute}.


\subsection{Model descriptions}
\label{sec:models}
% Scope of this subsection: what each model is and which weights we used.
% Sampling settings and prompts -> Hyperparameters. Serving software (vLLM),
% chat-template handling -> Experimental setup and code. GPUs, memory and
% runtimes -> Computational requirements.

We evaluate four open-weight large language models (LLMs). Two of them,
Mixtral-8x7B-Instruct-v0.1 and Llama-3-70B-Instruct, are the open models of the
original paper. We add two newer models, Qwen3.5-27B and DeepSeek-R1 7B. All
four are autoregressive Transformer decoders. Each one was pretrained on a
large text corpus and then post-trained to follow instructions. We do not
train or fine-tune any model. We use the published weights and classify
only through prompts. Table~\ref{tab:models} summarises the models.

\par\medskip

To support reproducibility, we record the model revisions and weight
precision used in our experiments. \addhere{Idk if this should be mentioned here} ** The original paper does not specify
the numerical precision or whether the weights were quantized. We contacted
the original authors for clarification, and they indicated that FP16 was
most likely used, although they could not confirm this with certainty.** 

\par\medskip

Accordingly, for Mixtral-8x7B-Instruct-v0.1, we use the published weights
in FP16 precision, without additional quantization.

\par\medskip

For Llama-3-70B-Instruct, the original precision is uncertain. The authors recall using INT8 but could not confirm it, and neither the method nor the checkpoint was recorded. Our initial INT8 attempts failed on the available hardware (Section~\ref{sec:setup}). We therefore use the checkpoint: RedHatAI/Meta-Llama-3-70B-Instruct-FP8, which is an FP8 quantization of Meta's weights published by a third party.

\par\medskip

For Qwen3.5-27B, we use the full, unquantized weights and fix the exact
revision of its Hugging Face repository. This identifies the weights,
tokenizer and chat template used. 

\par\medskip

For DeepSeek-R1 7B, we use the 4-bit
quantized weights (Q4\_K\_M) from the Ollama library and record their
model ID. A job stops if the model tag points to different weights.



\begin{table}[H]
  \centering
  \small
  \caption{The models we evaluate. Parameter counts and context lengths are
  from the model cards and the Mixtral paper. ``Active'' is the number of
  parameters used for each token.}
  \label{tab:models}
  % TODO: check the DeepSeek column against the model card before submission.
  \setlength{\tabcolsep}{4pt}
  \begin{tabular}{lcccc}
    \toprule
     & Mixtral-8x7B & Llama-3-70B & Qwen3.5-27B & DeepSeek-R1 7B \\
    \midrule
    Developer            & Mistral AI  & Meta              & Alibaba              & DeepSeek \\
    Architecture         & Sparse MoE  & Dense             & Dense, hybrid attn.  & Dense (distilled) \\
    Parameters (total)   & 47B         & 70B               & 27B                  & 7.6B \\
    Parameters (active)  & 13B         & 70B               & 27B                  & 7.6B \\
    Context length       & 32k         & 8k                & 262k                 & 131k \\
    Released             & Dec.\ 2023  & Apr.\ 2024        & Feb.\ 2026           & Jan.\ 2025 \\
    In original paper    & Yes         & Yes               & No                   & No \\
    Reasoning before answer & No       & No                & Off                  & On \\

    Revision used        & \texttt{eba9230} & \texttt{048a5cc}$^{*}$ & \texttt{fc05dae} & \texttt{755ced02ce7b}$^{**}$ \\
    
    Precision            & FP16             & FP8$^{*}$ & bfloat16         & 4-bit (Q4\_K\_M) \\

    License              & Apache 2.0  & Llama 3 Community & Apache 2.0           & MIT \\
    \bottomrule
  \end{tabular}

  \vspace{2pt}
  \raggedright\footnotesize
  $^{*}$ Third-party FP8 checkpoint \texttt{RedHatAI/Meta-Llama-3-70B-Instruct-FP8}.
  \\$^{**}$ ID of the tag \texttt{deepseek-r1:7b} in the Ollama library, not a Hugging Face revision.
\end{table}

\paragraph{Mixtral-8x7B-Instruct-v0.1~\cite{mixtral-8x7b-instruct-hf}.}

Mixtral-8x7B-Instruct-v0.1 is an instruction-tuned
language model developed by Mistral AI and released
in December 2023. It uses a sparse mixture-of-experts
(MoE) architecture with eight feed-forward experts
per layer, of which two are activated for each token.
The model contains approximately 47 billion parameters,
with around 13 billion active per token.

Mixtral is one of the models evaluated in the original
paper and is our primary reproduction target. We use
the published instruction-tuned weights from Hugging
Face (revision \texttt{eba9230}) in FP16 precision.

\paragraph{Llama-3-70B-Instruct~\cite{llama3-70b-instruct-hf}.}
Llama-3-70B-Instruct is an instruction-tuned language model from Meta,
released in April 2024. It is a dense decoder-only Transformer with 70 billion parameters and an 8k-token context window, so all parameters are used for each token. It is one of the two models of the original paper. It is about as old as Mixtral, but dense and larger.

\paragraph{Qwen3.5-27B~\cite{qwen35-27b-hf}.}
Qwen3.5-27B is a 27B-parameter model from the Qwen Team at Alibaba, released in
February 2026. It is about two years newer than the models of the original
paper. Its 64 layers mix two types of attention. Three of every four layers use
Gated DeltaNet, a form of linear attention. The fourth layer uses standard
gated attention. Every layer has a dense feed-forward block, so the model is
not a mixture of experts. The model also has a vision encoder, but our tasks
use only text input. The release contains the post-trained model. It has a
``thinking'' mode, on by default, that writes a reasoning trace before the
answer. We turn this mode off so that the model gives a direct answer, as
Mixtral and Llama do. We chose this model because its size is between the two
original models and it fits on the GPUs available to us. We use revision
\texttt{fc05dae} and bfloat16 weights.

\paragraph{DeepSeek-R1 7B~\cite{deepseek-r1-7b-ollama}.}
% TODO: add the bib entry deepseek-r1-7b-ollama.
DeepSeek-R1 7B is the smallest model in our study. It is
DeepSeek-R1-Distill-Qwen-7B, a distilled model: DeepSeek fine-tuned the
7.6B-parameter Qwen2.5-Math-7B on outputs of its large reasoning model
DeepSeek-R1. It was released in January 2025 under the MIT license.
We use the build \texttt{deepseek-r1:7b} of the Ollama library (ID
\texttt{755ced02ce7b}, 4.7\,GB, 4-bit quantization Q4\_K\_M). It differs from
the other three models in two ways. First, it is a reasoning model. It writes
about 250 tokens of reasoning before each answer. We tried to turn this off
(\texttt{reasoning\_effort=none} and \texttt{think=false}, on both the
OpenAI-compatible API and Ollama's own API), but the model reasoned in every
case, so we keep this default behaviour. Second, its weights are quantized. We chose this model to test the
two metrics on a model that is about ten times smaller than Llama-3-70B and
that runs on a single small GPU. Because of the two differences, its results
are not a pure test of model size (Section~\ref{sec:results}).
We ran DeepSeek-R1 7B on TREC and CB only. The reasoning makes it slow: at
about 1.4 requests per second on its node, RTE, DBPedia and WoS would each have
taken 36 to 45 hours, and we did not have this time
(Section~\ref{sec:compute}).

\subsection{Datasets}
\label{sec:datasets}
% Scope of this subsection: which data we evaluate, where it comes from, the
% splits, the sampling and the label changes. Prompt texts -> Experimental setup.
% Metrics and invalid answers (N/A) -> Experimental setup / Results.
% Problems we met while checking the data -> Discussion (What was difficult).

We use the five classification datasets of the original
paper~\cite{Errica:2025}: TREC, CommitmentBank (CB), RTE, DBPedia and Web of
Science (WoS). They cover question classification, natural language inference
(NLI) and topic classification, with 2 to 14 classes. We take the data files
from the authors' repository at commit \texttt{6573556}. We compared every file
with its official source (Table~\ref{tab:datasets}). The files are identical to
the official files, except for the small format changes that we list below.
No model is trained, so we use no training split for learning. Training or
validation files are used only as a source of the fixed one-shot examples.

\begin{table}[H]
  \centering
  \small
  \caption{The evaluated data. ``File'' is the official split that the
  authors' code reads. ``Evaluated'' is the number of samples that each model
  classifies. ``1-shot source'' is where the examples of the one-shot prompt
  come from.}
  \label{tab:datasets}
  \begin{tabular}{llcllr}
    \toprule
    Dataset & Task & Classes & File & 1-shot source & Evaluated \\
    \midrule
    TREC    & Question type & 6  & test (TREC-10)     & train            & 500  \\
    CB      & NLI           & 3  & train$^{\dagger}$  & validation       & 250  \\
    RTE     & NLI           & 2  & train$^{\ddagger}$ & fixed in code    & 2490 \\
    DBPedia & Topic         & 14 & test (70{,}000)    & unused test rows & 2000 \\
    WoS     & Topic         & 7  & all (46{,}985)     & unused rows      & 2000 \\
    \bottomrule
  \end{tabular}

  \vspace{2pt}
  \raggedright\footnotesize
  $^{\dagger}$ The paper calls these 250 samples ``test'' samples. We also evaluate Qwen3.5-27B and DeepSeek-R1 7B on the official test samples with rebuilt labels (see CB below).\\
  $^{\ddagger}$ The paper calls these 2490 samples ``test'' samples (see RTE below).
\end{table}

Figure~\ref{fig:class-counts} shows the class distribution of the evaluated
samples. TREC, CB and RTE use all rows of the file. DBPedia and WoS use a
balanced subset of 2000 samples. Two datasets are clearly imbalanced. In TREC,
the largest class is 15 times larger than the smallest class (Abbreviation, 9
samples). In CB, the Neutral class has only 16 samples (6.4\%). Per-class
results for these two classes rest on few samples. The figure also gives the
share of the majority class. This is the accuracy of a model that always
answers with the most common class, so it is the baseline for the accuracy of
each dataset: from 7.2\% for DBPedia to 50.2\% for RTE. For WoS, the full
dataset is imbalanced (Medical Science is 31\% of the 46{,}985 abstracts). The
authors' sampling removes this imbalance.

\begin{figure}[H]
  \centering
  % Made by figures/make_class_distribution.py (also writes an SVG copy).
  \includegraphics[width=\linewidth]{figures/class_distribution.pdf}
  \caption{Class distribution of the evaluated samples. Bars show the share of
  each class. Numbers at the bar ends are sample counts. The dashed line is the
  share that each class would have if the dataset were balanced. ``Majority''
  is the share of the largest class, and ``max/min'' is the size ratio of the
  largest to the smallest class. For WoS, the outline bars show the full
  dataset before the authors' balanced sampling.}
  \label{fig:class-counts}
\end{figure}

Table~\ref{tab:examples} shows one evaluated sample from each dataset. The
texts are copied from the data files without changes, including their spacing
and errors. The model sees the input text and must return the class. The CB
and RTE inputs have two parts: a premise and a hypothesis. The model must say
whether the premise implies the hypothesis.

\begin{table}[H]
  \centering
  \small
  \caption{One evaluated sample from each dataset, with its gold class.}
  \label{tab:examples}
  % Sources: TREC test.txt line 1; CB train.jsonl idx 146; RTE train.jsonl
  % idx 231; DBPedia evaluation subset id 37252; WoS evaluation subset id 26687.
  \begin{tabular}{lp{0.62\linewidth}l}
    \toprule
    Dataset & Input & Class \\
    \midrule
    TREC    & How far is it from Denver to Aspen ? & Number \\
    \addlinespace
    CB      & \textit{Premise:} A: Do you go to museums in Europe? B: Uh,
              actually, no, I don't think I went to any of them.\newline
              \textit{Hypothesis:} she went to some of them
            & Contradiction \\
    \addlinespace
    RTE     & \textit{Premise:} Yemen, too, was reunified in 1990.\newline
              \textit{Hypothesis:} Yemen was reunited in 1990.
            & Entailment \\
    \addlinespace
    DBPedia & Stillach is a river of Bavaria Germany. & Natural Place \\
    \addlinespace
    WoS     & Our paper studies design of microwave external applicator based
              on impedance matching. In this paper we deal with a topology and
              signal flow graph of applicator coupled to biological tissue.
            & Electrical Eng. \\
    \bottomrule
  \end{tabular}
\end{table}

\paragraph{TREC.}
TREC is a question classification dataset~\cite{Li:2002}. Each question has
one of six coarse classes, for example Number or Location. We evaluate all 500
questions of the TREC-10 test set~\cite{Voorhees:2000}. The six one-shot
examples come from the training set of 5452 questions, and none of them is in
the test set. Our test file has the same questions in the same order as the
official file. For example, Entity samples 46 to 50 are the five questions that
the paper quotes.

\paragraph{CommitmentBank (CB).}
CB is a three-class NLI dataset (entailment, contradiction,
neutral)~\cite{deMarneffe:2019}. We use the SuperGLUE version~\cite{Wang:2019},
which has 250 training, 56 validation and 250 test samples. The paper reports
250 test samples. However, the SuperGLUE test file has no labels, and the
authors' CB notebook reads the labelled training file, which also has 250
samples (119 contradiction, 115 entailment, 16 neutral). We follow the code
and evaluate all models on the 250 training samples. The three one-shot
examples come from the validation file, one per class. The training and test
files share no premises.

To also follow the paper's wording, we ran Qwen3.5-27B and DeepSeek-R1 7B a
second time on the official test samples. To get their labels, we matched each test sample with
its item in the original CB annotations and applied the published labelling
rule~\cite{Jiang:2019}: a label is given when at least 80\% of the annotator
ratings fall in one range. The same rule reproduces all 306 official training
and validation labels. The rebuilt test labels (121 contradiction, 113
entailment, 16 neutral) are a reconstruction, not an official release. We
report these runs separately (Sections~\ref{sec:results-qwen}
and~\ref{sec:results-deepseek}).

\paragraph{RTE.}
RTE is a binary NLI dataset from the PASCAL Recognising Textual Entailment
challenges~\cite{Dagan:2006}, in the SuperGLUE version~\cite{Wang:2019}.
SuperGLUE RTE has 2490 training, 277 validation and 3000 unlabelled test
samples. The paper reports 2490 test samples, and the authors' code reads the
training file. So we evaluate the 2490 training samples. No 2490-sample test
set exists. The code renames the class \texttt{not\_entailment} to
``contradiction/neutral''. The two one-shot examples are fixed in the code. One
of them is a CB validation item, not an RTE item. We keep both examples
unchanged.

\paragraph{DBPedia.}
DBPedia is a 14-class topic classification dataset built from Wikipedia
articles~\cite{Zhang:2015}. The model input is the article text, without its
title. The official test set has 70{,}000 samples. The authors' code shuffles
the test set with random seed 42 and keeps the first 142 ($= \lfloor 2000/14
\rfloor$) samples of each class. It then adds 12 random samples, which gives
2000 samples. The first unused sample of each class becomes its one-shot
example.

\paragraph{Web of Science (WoS).}
WoS contains abstracts of scientific papers~\cite{Kowsari:2017}. We use the
WOS-46985 version with its seven top-level classes, for example Computer
Science and Psychology~\cite{Kowsari:2018}. The full dataset is not balanced:
Medical Science has 14{,}625 abstracts, but Mechanical Engineering has only
3297. The authors' code uses the same procedure as for DBPedia, with
285 ($= \lfloor 2000/7 \rfloor$) samples for each class and 5 extra random
samples.

\paragraph{Preprocessing.}

We do not clean, tokenize or truncate any text. The prompts contain the
samples as they are in the files. The data files differ from the official files
only in format: the DBPedia file has an added header line, and one line of the
TREC training file has a different character encoding. That line is not
evaluated. We ran the authors' sampling code with Python 3.10 to 3.14. Every
version selected the same DBPedia and WoS samples, in the same order. These are
also the samples that the authors' stored notebook outputs show.

\paragraph{Availability.}

All datasets are public. TREC:
\url{https://cogcomp.seas.upenn.edu/Data/QA/QC/}. CB and RTE (SuperGLUE):
\url{https://super.gluebenchmark.com/tasks}. Original CB annotations:
\url{https://github.com/mcdm/CommitmentBank}. DBPedia:
\url{https://s3.amazonaws.com/fast-ai-nlp/dbpedia_csv.tgz}. WoS: Mendeley Data,
\url{https://doi.org/10.17632/9rw3vkcfy4.6}. Our repository contains the data
files, their SHA-256 hashes, the CB label script and the sampled subsets.

\subsection{Hyperparameters}
\label{sec:hparams}
No hyperparameters are trained or searched. Every value is fixed by the paper
or the released notebooks (Table~\ref{tab:hparams}); the only variable is the
model. Each model is evaluated in 15 configurations (5 datasets $\times$ 3
prompting strategies).

\par\medskip

The original paper and the publicly available setup
instructions do not specify the inference context
length used for Mixtral. In our reproduction, we
configured vLLM with a maximum context length of
2048 tokens for CB, RTE, and TREC, 3072 for
DBPedia, and 4096 for WoS. The larger limits
accommodated longer few-shot prompts, with the
longest WoS prompt containing 3274 tokens. 


\begin{table}[H]
  \centering
  \small
  \caption{Evaluation settings. ``Paper'' means Section~4 of~\cite{Errica:2025};
  ``code'' means the released notebooks.}
  \label{tab:hparams}
  \setlength{\tabcolsep}{4pt}
  \begin{tabular}{lll}
    \toprule
    Setting & Value & Source \\
    \midrule
    Prompt variants per sample, $Q$ & 30 (original + 29 rephrasings) & Paper, code \\
    Answers per variant, $A$ & 1 & Code \\
    Temperature (rephrasing, answering) & 0, 0 & Paper, code \\
    Seed & 42 & Paper, code \\
    Prompting strategies & Simple, detail, 1-shot & Paper, code \\
    Mixtral context length & 2048 / 3072 / 4096 tokens & Ours \\
    Llama context length & 2048 / 3072 / 4096 tokens & Ours \\
    Qwen thinking mode & Off & Ours \\
    DeepSeek reasoning & On (model default) & Ours \\
    DeepSeek answer limit & 2048 tokens (reasoning + answer) & Ours \\
    Qwen and DeepSeek context length & 8192 tokens & Ours \\
    \bottomrule
  \end{tabular}
\end{table}


\subsection{Experimental setup and code}
\label{sec:setup}
Figure~\ref{fig:pipeline} shows the pipeline. For each dataset, the evaluated
model first rewrites the task description 29 times (temperature 0); with the
original, this gives $Q = 30$ prompt variants $\rho_1,\dots,\rho_{30}$. The
same variants are reused for all three prompting strategies. The model then
classifies every sample $x$ once with every variant. The authors' extractor
checks the class names in a fixed order and assigns the first one that occurs
anywhere in the answer (case-insensitive substring match), or an extra class
N/A if none occurs. The 30 predictions give the empirical
distribution $p_\tau(y \mid x)$.

\begin{figure}[H]
  \centering
  % Made by figures/make_pipeline.py.
  \includegraphics[width=\linewidth]{figures/pipeline.pdf}
  \caption{Evaluation pipeline for one model, dataset and prompting strategy.}
  \label{fig:pipeline}
\end{figure}

\paragraph{Metrics.} We report the three metrics of the
paper~\cite{Errica:2025}: \emph{sensitivity}
$S_\tau(x) = H\!\left(p_\tau(\cdot \mid x)\right) / \ln (C+1)$, the
normalised entropy of the predictions for one sample, averaged over samples
(0 = the prompt variant never changes the answer); \emph{consistency}
$C_y = 1 - \overline{\mathrm{TVD}}$, one minus the total variation distance
between $p_\tau(\cdot \mid x)$ and $p_\tau(\cdot \mid x')$, averaged over all
pairs $x, x'$ of class $y$ (1 = all samples of a class behave alike); and the
\emph{micro-F1} score over all prompt variants. $C$ is the number of classes;
the $+1$ is the N/A class. The paper does not say whether N/A counts in the
normalisation. We include it because this gives values closer to Table~1 of
the paper. For consistency, the pairs include $x = x'$, and the classes are
pooled with a weight equal to their number of pairs. Each sample has exactly
one gold class and N/A counts as a wrong answer, so micro-F1 equals the share
of correct answers over all $30 \times$ samples answers. Sensitivity needs no
labels; consistency and F1 do.

\paragraph{Serving.} 

Mixtral-8x7B-Instruct-v0.1 was served using vLLM
with an OpenAI-compatible API. Inference used
FP16 precision, tensor parallel size 1, the Triton
MoE backend, and a GPU memory utilization setting
of 0.8.

We used the authors' original evaluation code,
with checkpointing added to accommodate the
three-hour Slurm job limit. This allowed interrupted
experiments to resume without repeating completed
evaluations.

\par\medskip

Llama-3-70B-Instruct was served with vLLM and its OpenAI-compatible API on a single GPU (tensor parallel size 1) with FP8 weights, a maximum of 64
concurrent sequences, a GPU memory utilization of 0.80 and the same maximum
context lengths as for Mixtral. Every job started its own server and ran the authors' evaluation on top of it. The server start took about 8 minutes per job, mostly for loading the weights.

\par\medskip

Qwen3.5-27B runs behind a local vLLM 0.32.0
server~\cite{Kwon:2023} with an OpenAI-compatible API. The authors' code sends
the requests in parallel through Ray. Qwen3.5-27B runs in eager mode with at
most 16 concurrent sequences, the settings under which its server started
reliably. Note that temperature 0 and a fixed seed do not make GPU inference
bitwise deterministic, because results can depend on batching.

\par\medskip

DeepSeek-R1 7B runs behind Ollama 0.40.4, which also has an OpenAI-compatible
API. Each job starts its own Ollama server, so that our settings do not depend
on the server of the node. Ollama handles 16 requests in parallel, and the
authors' code processes 2 samples at a time. We fix the context length to 8192
tokens, because Ollama cuts a prompt that is too long without an error.

\paragraph{Reasoning text.} The authors' extractor searches the whole answer
for class names. The reasoning of DeepSeek-R1 often names several classes, so
it must not reach the extractor. Ollama 0.40.4 returns the reasoning in a
separate field, and we pass only the final answer to the authors' code. The
rephrased task descriptions are handled in the same way. If the model reaches
the limit of 2048 tokens inside its reasoning, the answer is empty and the
extractor assigns N/A. This happened for at most 0.5\% of the DeepSeek
answers (Table~\ref{tab:deepseek-results}). The same limit also affects the
rephrasing step: in a few cases DeepSeek spent all 2048 tokens on reasoning and
wrote no rephrased task description, so that prompt variant has an empty
instruction line (1 of 30 variants for TREC, 2 of 30 for CB). The rest of the
prompt is unchanged. Without these variants every metric changes by at most
0.01.

\paragraph{Code.} Our code, job files, pinned package versions and data
checksums are at \url{TODO-anonymous-repository-link}.
% TODO: add an anonymised link (double-blind review: no author names).

\subsection{Computational requirements}
\label{sec:compute}
Mixtral-8x7B-Instruct-v0.1 was evaluated on the ITU HPC
DGX1 node, equipped with an NVIDIA GB10 GPU.
Each Slurm job was allocated 10 CPU cores and
115\,GB of memory, with a maximum runtime of
three hours per job with this setup. 

\par\medskip

Llama-3-70B-Instruct ran on the same DGX1 node as Mixtral: one NVIDIA GB10
GPU with 119\,GiB of memory shared by CPU and GPU, and a limit of three hours per job. The FP8 weights take about 70\,GB, so the model fits on this single GPU. The evaluation took 14 Slurm jobs and 33.0 hours in total. Starting the server took 2.0 of these hours (about 8 minutes per job), so the inference itself took 31.0 hours.

\par\medskip

Qwen3.5-27B runs on one node of the ITU HPC cluster with 2 NVIDIA L40S GPUs
(48\,GB each), 8 CPU cores and 60\,GB of RAM, one Slurm job per dataset.

\par\medskip

DeepSeek-R1 7B runs on a different node: an ARM machine with one NVIDIA GB10
chip and 128\,GB of memory shared by CPU and GPU, of which each job requests
32\,GB. The model uses about 6\,GB of it. The CB run on the training file
took 6.3 hours, about 1.0 request per second.
\addhere{add the DeepSeek TREC and CB-test times from the Slurm logs (jobs 182120 and 182119)}

\par\medskip

 Table~\ref{tab:compute}
lists the cost. For each dataset, every model that ran it gets the same
requests: samples $\times$ 30 variants $\times$ 3 strategies, 651{,}600 for
all five datasets. 

\par\medskip

Mixtral-8x7B-Instruct-v0.1 was initially tested on
other HPC nodes, but encountered insufficient GPU
memory. We therefore used the DGX1 node, equipped
with a single NVIDIA GB10 GPU with sufficient shared
memory to run the model in FP16 without quantization.
Inference used tensor parallel size 1. The evaluation
required approximately 113.8 GPU hours across the
five datasets, with WoS taking the longest (52.5 hours),
followed by DBPedia (28 hours) and RTE (22 hours).
Due to the three-hour Slurm job limit associated
with the required memory allocation, we implemented
checkpointing to resume interrupted evaluations.

\par\medskip

Qwen3.5-27B took 32.2 GPU
hours, about 11.5 requests per second; a preliminary test run and restarts
added 2.0 GPU hours. WoS is the slowest dataset per request because its inputs
are the longest (median abstract: 1304 characters).


\begin{table}[H]
  \centering
  \small
  \caption{Compute cost per model and dataset. Wall-clock time is the cumulative Slurm
  execution time, including model loading and restarts, but excluding queue time.
  GPU hours $=$ wall-clock time $\times$ number of GPUs. ``--'': not run.
  $^{\S}$~DeepSeek-R1 7B on the CB training file; its extra run on the CB test
  file is not included.}
  \label{tab:compute}
  \footnotesize
  \setlength{\tabcolsep}{3pt}
  \begin{tabular}{lr rr rr rr rr}
    \toprule
     & & \multicolumn{2}{c}{Qwen3.5-27B} & \multicolumn{2}{c}{DeepSeek-R1 7B} & \multicolumn{2}{c}{Mixtral-8x7B} & \multicolumn{2}{c}{Llama-3-70B} \\
    \cmidrule(lr){3-4} \cmidrule(lr){5-6} \cmidrule(lr){7-8} \cmidrule(lr){9-10}
    Dataset & Requests & Wall (h) & GPU h & Wall (h) & GPU h & Wall (h) & GPU h & Wall (h) & GPU h \\
    \midrule
    TREC    & 45{,}000  & 0.87  & 1.73  & \tbd & \tbd & 5.55  & 5.55  &  0.62  & 0.62 \\
    CB      & 22{,}500  & 0.83  & 1.66  & 6.30$^{\S}$ & 6.30 & 5.75  & 5.75  &  0.97  & 0.97 \\
    RTE     & 224{,}100 & 3.98  & 7.97  & --   & --   & 22.00 & 22.00 &  6.74  & 6.74 \\
    DBPedia & 180{,}000 & 4.16  & 8.31  & --   & --   & 28.00 & 28.00 &  5.54  & 5.54 \\
    WoS     & 180{,}000 & 6.26  & 12.51 & --   & --   & 52.50 & 52.50 & 19.14 & 19.14 \\
    \midrule
    Total   & 651{,}600 & 16.10 & 32.20 & \tbd & \tbd & 113.80 & 113.80 & 33.02 & 33.02 \\
    \bottomrule
  \end{tabular}
\end{table}

% TODO: fill the DeepSeek TREC time and total. Qwen GPU column sums to 32.18 (table says 32.20).

\section{Results}
\label{sec:results}

% DRAFT (to polish): overview.
All three claims hold for every model we ran: sensitivity, consistency and
micro-F1 change with the prompting strategy and the dataset, and no strategy
gives the best micro-F1 on every dataset. The paper's numbers are reproduced
closely for Mixtral on WoS, DBPedia and TREC (overall MAE at most 0.015) and
less closely on CB and RTE (0.053 and 0.060). For Llama-3-70B, sensitivity and
consistency are reproduced closely (MAE 0.008 and 0.005), micro-F1 only
partly (0.056). The best strategy matches the paper in 12 of 15
dataset--metric comparisons for Mixtral and 10 of 15 for Llama.

\par\medskip

We report the sensitivity, consistency, and
micro-F1 scores obtained across five classification
datasets and three prompting strategies.
We first compare the results for Mixtral-8x7B
and Llama-3-70B with those reported in the
original paper~\cite{Errica:2025}.
We then present the results for Qwen3.5-27B
and DeepSeek-R1 7B, which extend the original
evaluation to additional models.


\subsection{Results reproducing original paper}

\subsubsection{Mixtral-8x7B-Instruct-v0.1}

\paragraph{Reproduction of the original results.}
\mbox{}\par

Table~\ref{tab:mixtral-results} compares our
reproduced Mixtral results with the original
paper, while Table~\ref{tab:mixtral-mae}
summarizes the mean absolute error (MAE)
for each dataset and metric.

\par\medskip

MAE was calculated as the average absolute
difference between our results and the published
values across the three prompting strategies.
The overall MAE for each dataset was obtained
by averaging across all three metrics and
prompting strategies.

\par\medskip

The closest numerical agreement was observed
for WoS (overall MAE: 0.0053), followed by
DBPedia (0.0134) and TREC (0.0147).
Larger differences occurred for CB (0.0529)
and RTE (0.0604). RTE also showed the largest
sensitivity MAE (0.0519) and the largest
micro-F1 MAE (0.0793). For CB we use corrected
reference values for two sensitivity cells
(note of Table~\ref{tab:mixtral-results}); with
the values printed in the paper, the CB
sensitivity MAE would be 0.1281.

\begin{table}[H]
\centering
\caption{Mixtral-8x7B reproduction results compared
with the original paper. $\Delta$ = Ours $-$ Paper.}
\label{tab:mixtral-results}
\scriptsize
\setlength{\tabcolsep}{3pt}
\resizebox{\textwidth}{!}{%
\begin{tabular}{llrrrrrrrrr}
\toprule
& & \multicolumn{3}{c}{Sensitivity}
    & \multicolumn{3}{c}{Consistency}
    & \multicolumn{3}{c}{Micro-F1} \\
\cmidrule(lr){3-5}
\cmidrule(lr){6-8}
\cmidrule(lr){9-11}
Dataset & Strategy
& Ours & Paper & $\Delta$
& Ours & Paper & $\Delta$
& Ours & Paper & $\Delta$ \\
\midrule
TREC & Simple
& 0.1906 & 0.223 & -0.0324
& 0.6566 & 0.657 & -0.0004
& 0.7294 & 0.736 & -0.0066 \\
& Detail
& 0.1919 & 0.182 & +0.0099
& 0.6834 & 0.709 & -0.0256
& 0.7056 & 0.728 & -0.0224 \\
& 1-shot
& 0.1429 & 0.166 & -0.0231
& 0.6463 & 0.640 & +0.0063
& 0.7380 & 0.744 & -0.0060 \\
\midrule
CB & Simple
& 0.2531 & 0.203 & +0.0501
& 0.5607 & 0.625 & -0.0643
& 0.6001 & 0.672 & -0.0719 \\
& Detail
& 0.1794 & 0.125$^{\ast}$ & +0.0544
& 0.6215 & 0.695 & -0.0735
& 0.7032 & 0.788 & -0.0848 \\
& 1-shot
& 0.1847 & 0.174$^{\ast}$ & +0.0107
& 0.5591 & 0.600 & -0.0409
& 0.6263 & 0.652 & -0.0257 \\
\midrule
RTE & Simple
& 0.2599 & 0.195 & +0.0649
& 0.5945 & 0.669 & -0.0745
& 0.6931 & 0.814 & -0.1209 \\
& Detail
& 0.0994 & 0.176 & -0.0766
& 0.7511 & 0.676 & +0.0751
& 0.8501 & 0.756 & +0.0941 \\
& 1-shot
& 0.0601 & 0.046 & +0.0141
& 0.7545 & 0.754 & +0.0005
& 0.8241 & 0.847 & -0.0229 \\
\midrule
DBPedia & Simple
& 0.0494 & 0.053 & -0.0036
& 0.8155 & 0.802 & +0.0135
& 0.8649 & 0.804 & +0.0609 \\
& Detail
& 0.0537 & 0.054 & -0.0003
& 0.8112 & 0.808 & +0.0032
& 0.8681 & 0.862 & +0.0061 \\
& 1-shot
& 0.0728 & 0.069 & +0.0038
& 0.7682 & 0.773 & -0.0048
& 0.8432 & 0.819 & +0.0242 \\
\midrule
WoS & Simple
& 0.1117 & 0.112 & -0.0003
& 0.4845 & 0.483 & +0.0015
& 0.5945 & 0.606 & -0.0115 \\
& Detail
& 0.1284 & 0.124 & +0.0044
& 0.4682 & 0.469 & -0.0008
& 0.6020 & 0.599 & +0.0030 \\
& 1-shot
& 0.1475 & 0.142 & +0.0055
& 0.4756 & 0.472 & +0.0036
& 0.5492 & 0.566 & -0.0168 \\
\bottomrule
\end{tabular}%
}

\vspace{2pt}
\raggedright\footnotesize
$^{\ast}$ Table~1 of the paper gives 0.013 and 0.017. We use 0.125 and 0.174,
the values in the authors' stored notebook outputs. Appendix~D of the paper
gives a standard deviation of 0.188 for both cells, which is not possible for a
mean of 0.013 on a $[0,1]$ scale, so Table~1 most likely has a decimal error.
With the printed values, $\Delta$ would be $+0.166$ and $+0.168$.
\end{table}

\begin{table}[H]
\centering
\caption{Mean absolute error (MAE) between reproduced
and published Mixtral results, averaged across
the three prompting strategies for each metric.}
\label{tab:mixtral-mae}
\small
\begin{tabular}{lrrrr}
\toprule
Dataset & Sensitivity & Consistency & Micro-F1 & Overall MAE \\
\midrule
TREC    & 0.0218 & 0.0108 & 0.0117 & 0.0147 \\
CB      & 0.0384 & 0.0596 & 0.0608 & 0.0529 \\
RTE     & 0.0519 & 0.0500 & 0.0793 & 0.0604 \\
DBPedia & 0.0026 & 0.0072 & 0.0304 & 0.0134 \\
WoS     & 0.0034 & 0.0020 & 0.0104 & 0.0053 \\
\bottomrule
\end{tabular}
\end{table}

\paragraph{Reproduction of the original findings.}
\mbox{}\par

Our Mixtral results support the three claims
outlined in Section~\ref{sec:claims}.
Sensitivity (Claim 1), consistency (Claim 2),
and micro-F1 (Claim 3) vary across prompting
strategies and datasets. No single prompting
strategy consistently achieves the best results
across all datasets and metrics, in agreement
with the general findings of the original
paper~\cite{Errica:2025}.

\par\medskip

The best-performing prompting strategy matched
the original paper in 12 of 15 dataset--metric
comparisons (80\%), including all five
sensitivity comparisons. The three disagreements
concerned micro-F1 on WoS and RTE, and
consistency on DBPedia.

\par\medskip

Overall, the Mixtral experiment partially
reproduced the original findings. Although
some numerical differences remained, most
of the best-performing prompting strategies
were reproduced.


\subsubsection{Llama-3-70B-Instruct}

Table~\ref{tab:llama-results} compares our Llama-3-70B-Instruct results with the original paper, and Table~\ref{tab:llama-mae} gives the mean absolute error (MAE), computed in the same way as for Mixtral. Micro-F1 is the average over the 30 prompt variants. ``Detail'' and ``1-shot'' are the paper's names for the \texttt{instruct} and \texttt{fewshot} prompts.

Sensitivity and consistency are reproduced closely. The MAE is 0.0081 for
sensitivity and 0.0048 for consistency. The largest differences are 0.023
(sensitivity, DBPedia with the detailed prompt) and 0.015 (consistency, TREC with the one-shot prompt). Micro-F1 differs more (MAE 0.0564). The closest agreement is on CB (overall MAE 0.0045) and WoS (0.0056), followed by RTE (0.0264), TREC (0.0321) and DBPedia (0.0467).

The F1 differences are concentrated in a few configurations. On TREC, our F1 is lower than the paper's by 0.06 to 0.09 for all three prompts. On RTE with the detailed prompt, our F1 is 0.161 higher (0.670 against 0.509). On DBPedia, it is 0.206 higher with the detailed prompt (0.946 against 0.740) and 0.134 higher with the one-shot prompt (0.961 against 0.827). The paper's values for these two DBPedia configurations are much lower than its own value for the simple prompt (0.948), while ours stay between 0.946 and 0.965.\addhere{Add the share of N/A answers for these configurations, from the evaluation script.}

\begin{table}[H]
\centering
\caption{Llama-3-70B-Instruct reproduction results compared with the original paper. $\Delta$ = Ours $-$ Paper. Our model is the FP8 checkpoint of RedHat AI.}
\label{tab:llama-results}
\scriptsize
\setlength{\tabcolsep}{3pt}
\resizebox{\textwidth}{!}{%
\begin{tabular}{llrrrrrrrrr}
\toprule
& & \multicolumn{3}{c}{Sensitivity}
    & \multicolumn{3}{c}{Consistency}
    & \multicolumn{3}{c}{Micro-F1} \\
\cmidrule(lr){3-5}
\cmidrule(lr){6-8}
\cmidrule(lr){9-11}
Dataset & Strategy
& Ours & Paper & $\Delta$
& Ours & Paper & $\Delta$
& Ours & Paper & $\Delta$ \\
\midrule
TREC & Simple
& 0.1376 & 0.127 & +0.0106
& 0.6923 & 0.693 & -0.0007
& 0.7558 & 0.848 & -0.0922 \\
& Detail
& 0.1605 & 0.141 & +0.0195
& 0.6824 & 0.694 & -0.0116
& 0.7480 & 0.824 & -0.0760 \\
& 1-shot
& 0.0906 & 0.095 & -0.0044
& 0.7005 & 0.686 & +0.0145
& 0.7865 & 0.846 & -0.0595 \\
\midrule
CB & Simple
& 0.0254 & 0.018 & +0.0074
& 0.9551 & 0.956 & -0.0009
& 0.9297 & 0.920 & +0.0097 \\
& Detail
& 0.0186 & 0.016 & +0.0026
& 0.9630 & 0.963 & 0.0000
& 0.9313 & 0.924 & +0.0073 \\
& 1-shot
& 0.0186 & 0.014 & +0.0046
& 0.9557 & 0.962 & -0.0063
& 0.9345 & 0.936 & -0.0015 \\
\midrule
RTE & Simple
& 0.1591 & 0.157 & +0.0021
& 0.6994 & 0.701 & -0.0016
& 0.6123 & 0.569 & +0.0433 \\
& Detail
& 0.2133 & 0.213 & +0.0003
& 0.7268 & 0.727 & -0.0002
& 0.6698 & 0.509 & +0.1608 \\
& 1-shot
& 0.0523 & 0.048 & +0.0043
& 0.7497 & 0.755 & -0.0053
& 0.7986 & 0.818 & -0.0194 \\
\midrule
DBPedia & Simple
& 0.0137 & 0.018 & -0.0043
& 0.9360 & 0.932 & +0.0040
& 0.9651 & 0.948 & +0.0171 \\
& Detail
& 0.0350 & 0.058 & -0.0230
& 0.9158 & 0.901 & +0.0148
& 0.9455 & 0.740 & +0.2055 \\
& 1-shot
& 0.0180 & 0.030 & -0.0120
& 0.9301 & 0.924 & +0.0061
& 0.9609 & 0.827 & +0.1339 \\
\midrule
WoS & Simple
& 0.0476 & 0.043 & +0.0046
& 0.4765 & 0.479 & -0.0025
& 0.6288 & 0.627 & +0.0018 \\
& Detail
& 0.0405 & 0.034 & +0.0065
& 0.4812 & 0.480 & +0.0012
& 0.6382 & 0.623 & +0.0152 \\
& 1-shot
& 0.0446 & 0.030 & +0.0146
& 0.4696 & 0.468 & +0.0016
& 0.5825 & 0.580 & +0.0025 \\
\bottomrule
\end{tabular}%
}
\end{table}

\begin{table}[H]
\centering
\caption{Mean absolute error (MAE) between reproduced and published
Llama-3-70B-Instruct results, averaged across the three prompting strategies
for each metric.}
\label{tab:llama-mae}
\small
\begin{tabular}{lrrrr}
\toprule
Dataset & Sensitivity & Consistency & Micro-F1 & Overall MAE \\
\midrule
TREC    & 0.0115 & 0.0089 & 0.0759 & 0.0321 \\
CB      & 0.0049 & 0.0024 & 0.0062 & 0.0045 \\
RTE     & 0.0022 & 0.0024 & 0.0745 & 0.0264 \\
DBPedia & 0.0131 & 0.0083 & 0.1188 & 0.0467 \\
WoS     & 0.0086 & 0.0018 & 0.0065 & 0.0056 \\
\midrule
All     & 0.0081 & 0.0048 & 0.0564 & 0.0231 \\
\bottomrule
\end{tabular}
\end{table}

\paragraph{Reproduction of the original findings.}
\mbox{}\par

Our Llama results support the three claims of Section~\ref{sec:claims}.
Sensitivity (Claim 1), consistency (Claim 2) and micro-F1 (Claim 3) change
with the prompting strategy and the dataset. For example, the sensitivity on RTE is 0.213 with the detailed prompt and 0.052 with the one-shot prompt. No strategy gives the best micro-F1 on every dataset: it is the one-shot prompt on TREC, CB and RTE, the simple prompt on DBPedia and the detailed prompt on WoS.

The best strategy of our run matches the paper in 10 of the 15
dataset--metric comparisons. It differs in 4: sensitivity and micro-F1 on WoS, and consistency and micro-F1 on TREC. The remaining case, CB sensitivity, is a tie: the detailed and the one-shot prompt both have 0.0186 in our run, whereas the paper has the one-shot prompt as best. The differences on WoS and on TREC consistency are small (at most 0.015 between strategies in our run). Overall, the Llama experiment reproduces sensitivity and consistency well and reproduces micro-F1 only partly. We did not test whether FP8 weights, which differ from the INT8 weights the authors recall, explain the F1 differences.

\subsection{Results beyond original paper}

\subsubsection{Qwen3.5-27B}
\label{sec:results-qwen}

Table~\ref{tab:qwen-results} reports the Qwen3.5-27B results. The metrics are
computed with the same code as for Mixtral. The main CB rows use the training
file, so they cover the same 250 samples as our Mixtral run. The extra CB rows
use the official test samples with our rebuilt labels
(Section~\ref{sec:datasets}). N/A is the share of the 30 answers per sample
for which the extractor found no class.

\begin{table}[H]
\centering
\caption{Qwen3.5-27B results. CB (train) is the main CB run, on the same
samples as Mixtral. CB (test) uses the official test samples with rebuilt
labels and is not comparable with the paper's Table~1.}
\label{tab:qwen-results}
\small
% Computed from qwen3.5-27b-results.zip in the code repository (commit 1dd67f8),
% with the formulas of mixtral8x7b-Instruct-v0.1_RESULTS/Results_and_paper_comparison.ipynb.
% CB (train) = Slurm job 182903.
\begin{tabular}{llrrrr}
\toprule
Dataset & Strategy & Sensitivity & Consistency & Micro-F1 & N/A (\%) \\
\midrule
TREC & Simple  & 0.0344 & 0.7539 & 0.8427 & 0.00 \\
     & Detail  & 0.0361 & 0.7292 & 0.8273 & 0.03 \\
     & 1-shot  & 0.0288 & 0.7892 & 0.8598 & 0.00 \\
\midrule
CB (train) & Simple  & 0.0081 & 0.8736 & 0.9280 & 0.00 \\
           & Detail  & 0.0081 & 0.8736 & 0.9272 & 0.00 \\
           & 1-shot  & 0.0156 & 0.8562 & 0.9199 & 0.00 \\
\midrule
RTE  & Simple  & 0.0890 & 0.8039 & 0.8841 & 5.46 \\
     & Detail  & 0.1768 & 0.7853 & 0.8530 & 5.99 \\
     & 1-shot  & 0.0138 & 0.8886 & 0.9407 & 0.00 \\
\midrule
DBPedia & Simple  & 0.0026 & 0.9537 & 0.9746 & 0.00 \\
        & Detail  & 0.0030 & 0.9501 & 0.9724 & 0.00 \\
        & 1-shot  & 0.0021 & 0.9517 & 0.9734 & 0.00 \\
\midrule
WoS  & Simple  & 0.0303 & 0.4941 & 0.6588 & 0.26 \\
     & Detail  & 0.0253 & 0.4866 & 0.6480 & 0.07 \\
     & 1-shot  & 0.0223 & 0.4925 & 0.6510 & 0.05 \\
\midrule
CB (test) & Simple  & 0.0211 & 0.7938 & 0.8868 & 0.00 \\
          & Detail  & 0.0159 & 0.7693 & 0.8716 & 0.00 \\
          & 1-shot  & 0.0117 & 0.7604 & 0.8631 & 0.00 \\
\bottomrule
\end{tabular}
\end{table}

Sensitivity is below 0.04 in 13 of the 15 main configurations. The two
exceptions are RTE with the simple (0.089) and detailed (0.177) prompts, which
are also the only configurations with more than 1\% N/A answers. With the
one-shot prompt, RTE has no N/A answers and a sensitivity of 0.014. The best
micro-F1 comes from the one-shot prompt for TREC and RTE, and from the simple
prompt for CB, DBPedia and WoS.

On the CB training samples, Qwen3.5-27B reaches a micro-F1 of 0.920 to 0.928,
against 0.600 to 0.703 for Mixtral on the same samples
(Table~\ref{tab:mixtral-results}). Its sensitivity is 0.008 to 0.016, against
0.179 to 0.253 for Mixtral.

On the CB test samples, Qwen3.5-27B has a micro-F1 that is 0.041 to 0.057
lower and a consistency that is 0.080 to 0.104 lower than on the training
samples, for every prompt. Sensitivity stays below 0.022 on both sample sets.
The model and prompts are the same in both runs, so these differences come
from the samples and their labels.

\subsubsection{DeepSeek-R1 7B}
\label{sec:results-deepseek}

Table~\ref{tab:deepseek-results} reports the DeepSeek-R1 7B results for the
two datasets it ran on. The metrics are computed with the same code as for the
other models. As for Qwen3.5-27B, the main CB rows use the training file and
the extra CB rows the official test samples with our rebuilt labels.

\begin{table}[H]
\centering
\caption{DeepSeek-R1 7B results. CB (train) is the main CB run, on the same
samples as the other models. CB (test) uses the official test samples with
rebuilt labels. RTE, DBPedia and WoS were not run.}
\label{tab:deepseek-results}
\small
% Computed with scripts/deepseek/compare_models.py (branch deepseek-r1-7b,
% commit 6a17f70). TREC = Slurm job 182120, CB (test) = 182119, CB (train) = 183238.
\begin{tabular}{llrrrr}
\toprule
Dataset & Strategy & Sensitivity & Consistency & Micro-F1 & N/A (\%) \\
\midrule
TREC & Simple  & 0.1668 & 0.7353 & 0.7827 & 0.19 \\
     & Detail  & 0.1645 & 0.7392 & 0.7738 & 0.45 \\
     & 1-shot  & 0.1350 & 0.7634 & 0.8039 & 0.11 \\
\midrule
CB (train) & Simple  & 0.2769 & 0.6786 & 0.7399 & 0.21 \\
           & Detail  & 0.2994 & 0.6777 & 0.7325 & 0.25 \\
           & 1-shot  & 0.2971 & 0.6442 & 0.7039 & 0.27 \\
\midrule
CB (test) & Simple  & 0.2764 & 0.6669 & 0.7439 & 0.35 \\
          & Detail  & 0.2960 & 0.6699 & 0.7400 & 0.41 \\
          & 1-shot  & 0.3060 & 0.6443 & 0.7200 & 0.31 \\
\bottomrule
\end{tabular}
\end{table}

The claims hold for DeepSeek-R1 7B on both datasets: the metrics change with
the prompting strategy, and the best strategy depends on the dataset. On TREC
the one-shot prompt is best on all three metrics; on CB (train) the simple
prompt is.

On CB (train), DeepSeek-R1 7B is the most sensitive of the four models:
0.277 to 0.299, against 0.179 to 0.253 for Mixtral on the same samples. Yet its
micro-F1 is higher than Mixtral's for every prompt (0.704 to 0.740, against
0.600 to 0.703), and so is its consistency. So on CB, sensitivity and micro-F1
rank these two models in opposite order. On TREC, DeepSeek-R1 7B is better than
Mixtral on all three metrics.

DeepSeek-R1 7B gives almost the same results on the CB training and test
samples: no metric differs by more than 0.016. Both CB runs use the same 30
prompt variants, because the rephrasing step sees only the task description
and runs at temperature 0 with a fixed seed. This is unlike Qwen3.5-27B, whose
micro-F1 drops by 0.041 to 0.057 on the test samples
(Section~\ref{sec:results-qwen}).

\section{Discussion}

% DRAFT (to polish).
Our results support the three claims of the paper for all four models. The
paper's numbers are only partly reproducible: closely for most of Mixtral and
for the sensitivity and consistency of Llama-3-70B, less so for Mixtral on CB
and RTE and for the micro-F1 of Llama. We could not separate the possible
causes, because the paper does not record the serving setup and the weight
precision. Our main limits are one run per configuration, FP8 instead of the
recalled INT8 weights for Llama, and only two datasets for DeepSeek-R1 7B.

% DRAFT (to polish). Numbers: means over the three prompting strategies, computed
% with the same code as Tables~\ref{tab:mixtral-results} and \ref{tab:qwen-results}.
\subsection{Comparing different models}
\label{sec:disc-models}

The original paper uses two models from 2023--2024. We also ran Qwen3.5-27B,
which is about two years newer, and DeepSeek-R1 7B, a small reasoning model.
We have full results for Mixtral, Llama-3-70B and Qwen3.5-27B, and results
on TREC and CB for DeepSeek-R1 7B.

\begin{table}[H]
\centering
\caption{Mean over the three prompting strategies for each model and dataset.
CB uses the training file for all models. Bold: best value per dataset and
metric among the models with results. ``--'': not run. Per-strategy values
are in Tables~\ref{tab:mixtral-results}, \ref{tab:llama-results},
\ref{tab:qwen-results} and~\ref{tab:deepseek-results}.}
\label{tab:model-comparison}
\scriptsize
\setlength{\tabcolsep}{3pt}
\resizebox{\textwidth}{!}{%
\begin{tabular}{l rrrr rrrr rrrr}
\toprule
& \multicolumn{4}{c}{Sensitivity $\downarrow$}
& \multicolumn{4}{c}{Consistency $\uparrow$}
& \multicolumn{4}{c}{Micro-F1 $\uparrow$} \\
\cmidrule(lr){2-5} \cmidrule(lr){6-9} \cmidrule(lr){10-13}
Dataset & Mixtral & Llama & Qwen & DeepSeek
        & Mixtral & Llama & Qwen & DeepSeek
        & Mixtral & Llama & Qwen & DeepSeek \\
\midrule
TREC    & 0.175 & 0.130 & \textbf{0.033} & 0.155 & 0.662 & 0.692 & \textbf{0.757} & 0.746 & 0.724 & 0.763 & \textbf{0.843} & 0.787 \\
CB      & 0.206 & 0.021 & \textbf{0.011} & 0.291 & 0.580 & \textbf{0.958} & 0.868 & 0.667 & 0.643 & \textbf{0.932} & 0.925 & 0.725 \\
RTE     & 0.140 & 0.142 & \textbf{0.093} & --    & 0.700 & 0.725 & \textbf{0.826} & --    & 0.789 & 0.694 & \textbf{0.893} & --    \\
DBPedia & 0.059 & 0.022 & \textbf{0.003} & --    & 0.798 & 0.927 & \textbf{0.952} & --    & 0.859 & 0.957 & \textbf{0.973} & --    \\
WoS     & 0.129 & 0.044 & \textbf{0.026} & --    & 0.476 & 0.476 & \textbf{0.491} & --    & 0.582 & 0.617 & \textbf{0.653} & --    \\
\bottomrule
\end{tabular}%
}
\end{table}

\paragraph{What we observe.}
Qwen3.5-27B has lower sensitivity, higher consistency and higher micro-F1
than Mixtral on all five datasets (Table~\ref{tab:model-comparison}). The
biggest difference is on CB. Sensitivity goes from 0.206 to 0.011 and
micro-F1 from 0.643 to 0.925. The smallest differences are the sensitivity on
RTE (0.140 vs.\ 0.093) and the consistency on WoS (0.476 vs.\ 0.491).
Mixtral is most sensitive on CB, while Qwen3.5-27B is most sensitive on RTE.

The claims of the paper (Section~\ref{sec:claims}) hold for both models. The
three metrics change with the prompting strategy and the dataset, and no
strategy gives the best micro-F1 on every dataset. But how large the effect
is depends a lot on the model. Qwen3.5-27B has a sensitivity below 0.04 in 13
of the 15 configurations, so the rephrasings almost never change its answer.

\paragraph{Possible explanations.}
We did not run extra experiments to find out why the models differ. Below we
give some possible reasons, based on the results we have. They are
hypotheses, not tested findings. Section~\ref{sec:future-work} lists what
could be done to test them.

\paragraph{Each model writes its own rephrasings.}
In the authors' code, the model that is evaluated also writes the 29
rephrasings of the task description (Section~\ref{sec:setup}). So Mixtral and
Qwen3.5-27B do not get the same prompts. We compared the two sets. Mixtral's
rephrasings are longer (median 27 to 42 words, depending on the dataset,
against 17 to 32 for Qwen3.5-27B). On RTE, all 29 of Mixtral's rephrasings
are yes/no questions, and none of Qwen3.5-27B's are. Mixtral also sometimes
copies its own instructions into the rephrasing, for example ``Additionally,
avoid generating the same question as those listed below''. This happens in
12 of the 29 CB rephrasings and 6 of the 29 TREC rephrasings, and never for
Qwen3.5-27B, whose rephrasings are short instructions close to the original.
Prompts that are more alike should change the answers less, so part of the
lower sensitivity of Qwen3.5-27B may come from its rephrasings and not from
how it classifies. Other studies found that small changes in prompt format
can change accuracy by many points, differently for each
model~\cite{Sclar:2024}, and that model scores and rankings change across
paraphrases of the instruction~\cite{Mizrahi:2024}.

\paragraph{Model generation.}
Qwen3.5-27B is newer and probably better at following instructions. The
number of N/A answers (no class found in the answer) points in this
direction. It is below 1.5\% for both models on most datasets. The exception
is RTE with the simple and detailed prompts: 21.8\% and 0.9\% for Mixtral,
and 5.5\% and 6.0\% for Qwen3.5-27B. With only two models, we cannot tell if
the difference comes from the training, the size or the architecture
(Table~\ref{tab:models}). Other work found larger models to be more robust to
prompt changes~\cite{Zhuo:2024}, but also that size and instruction tuning do
not remove the sensitivity to prompt format~\cite{Sclar:2024}. 
Llama-3-70B is from the same time as Mixtral, but dense and larger. It has
lower sensitivity than Mixtral on four of the five datasets (for example CB: 0.021 against 0.206) and about the same on RTE (0.142 against 0.140). Its consistency is higher on four datasets and equal on WoS. Its micro-F1 is higher on four datasets and lower on RTE (0.694 against 0.789). Qwen3.5-27B is better than Llama on every dataset except CB, where Llama has the higher consistency (0.958 against 0.868) and micro-F1 (0.932 against 0.925). So a model of the same age as Mixtral is already much more stable on most datasets, which suggests that size or architecture matter and not only the model generation. With three models, and FP8 weights for Llama
against FP16 for Mixtral, we cannot separate these factors.
DeepSeek-R1 7B, the smallest model, does not fit a simple size order either.
On TREC it is less sensitive than Mixtral (0.155 against 0.175) and has higher
consistency and micro-F1 than Mixtral and Llama-3-70B. On CB it is the most
sensitive model (0.291), although its micro-F1 is higher than Mixtral's. It
also differs from the other models in more than size: it reasons before each
answer and runs with 4-bit weights.

\paragraph{RTE and the answer extractor.}
The RTE classes are ``entailment'' and ``contradiction/neutral'', and the
authors' extractor only accepts these two strings. If the model answers just
``contradiction'' or just ``neutral'', the answer becomes N/A. Qwen3.5-27B's
rephrasings of the RTE task mention three options (``entailment,
contradiction, or neutral''), which can lead to such answers. This could
explain why RTE is the only dataset with many N/A answers for Qwen3.5-27B,
and also its highest sensitivity, since N/A counts as an extra class. With
the one-shot prompt, where an example shows the label
``contradiction/neutral'', Qwen3.5-27B has no N/A answers and a sensitivity
of 0.014. We cannot check this directly, because the result files only keep
the extracted class and not the full answer. Answer extraction is known to
affect reported scores in other benchmarks too~\cite{Molfese:2025}.

\paragraph{Training data and the CB split.}
On CB, Qwen3.5-27B scores 0.041 to 0.057 lower micro-F1 on the test samples
than on the training samples (Section~\ref{sec:results-qwen}). One reason
could be that Qwen3.5-27B saw the CB training data during pretraining. Parts
of GLUE and SuperGLUE have been found in large pretraining
corpora~\cite{Elazar:2024}, but we do not know if CB is one of them or what
data Qwen3.5-27B was trained on. The gap can also simply come from the two
sample sets being different. DeepSeek-R1 7B speaks somewhat against this
second explanation: with the same prompt variants, it scores almost the same
on both sample sets (no metric differs by more than 0.016,
Section~\ref{sec:results-deepseek}). We do not think our rebuilt labels cause it,
since the same rule gives back all the official training labels.

\paragraph{Consistency on WoS.}
On WoS, consistency is almost the same for both models (0.476 and 0.491),
even though Qwen3.5-27B has higher micro-F1 and lower sensitivity. Both models
stay below 0.66 micro-F1. Consistency compares the answers for samples of the
same class, so it stays low when a model gives different wrong classes to
samples of one class. Some WoS topics are close to each other, for example
Computer Science and Electrical Engineering. This may limit consistency for
any model on this dataset.

\paragraph{How the models were run.}
Mixtral ran in FP16, Llama-3-70B in FP8 and Qwen3.5-27B in bfloat16, on
different GPUs and with different vLLM settings, and DeepSeek-R1 7B ran with
4-bit weights behind Ollama (Section~\ref{sec:setup}). Even at temperature 0, GPU
inference is not fully deterministic. We expect this to matter much less
than the differences above, but we did not test it.

\subsection{What was easy}
Several parts of the reproduction were easy because the authors' release is
well structured. The paper states the key settings (temperature 0, seed 42,
30 prompt variants, one answer per variant), and the code has one evaluation notebook per dataset and a table of model names. This let us add new models without changing the evaluation logic. The result files store the extracted class of all 30 variants for every sample. We could therefore recompute sensitivity, consistency and micro-F1 ourselves and check our scripts against the authors' evaluation functions on a small run. The sampling of DBPedia and WoS is seeded, and it selected the same samples with every Python version we tried, so the evaluated samples are well defined.

\par\medskip

Once the first dataset ran, the others were easy to add. We wrote the
runners for the other datasets by following the structure of the first one
and tested them on a CPU job without the model, which found small bugs
(for example the missing imports in the DBPedia runner) without waiting for
the model to load. Checkpointing made the three-hour job limit harmless,
because a resubmitted job skips the samples that are finished. These points
were easy for us because we had a working serving setup and the code needed
only small changes. Someone without access to a GPU node of this size would
find the compute much harder than the code.

\subsection{What was difficult}
% DRAFT (to polish): CB split decision. Facts in Section~\ref{sec:datasets}.
\paragraph{The CB split.}
Our aim was maximum reproducibility. So we checked every data file against its
official source before we ran any model (Section~\ref{sec:datasets}). For CB,
the paper and the released code do not agree. The paper reports 250
\emph{test} samples. The official SuperGLUE CB test file has 250 samples, but
it has no labels. The authors' code reads a label from every sample, so it
cannot run on that file. The code reads the training file, which also has 250
samples, with labels. Thus, the authors most likely evaluated the training
samples and called them test samples. We cannot prove this, because the paper
does not name the file and the row count alone does not identify it.

We had two options. Option A: use the training file, as the code does. This
probably gives the same samples as the paper's numbers. Option B: use the test
samples, as the paper text says. This needs labels that do not exist in the
official release. Our main runs use option A: the Mixtral reproduction ran the
authors' notebook unchanged on the training file, and we ran Qwen3.5-27B on
the same file, so both models are compared on the same samples. The same
reasoning applies to RTE, where the paper reports 2490 test samples but the
only file of that size is the training file that the code reads.

We also tried option B for Qwen3.5-27B and DeepSeek-R1 7B. We rebuilt the test labels from the
original CB annotations with the published labelling rule~\cite{Jiang:2019}:
a label is given when at least 80\% of the annotators agree. To check the
method, we applied the same rule to the training and validation files. It
gave back all 306 official labels. The two sample sets are of the same size
and have a similar class balance (training: 119 contradiction, 115
entailment, 16 neutral; test: 121, 113, 16), but they are not the same items.
So a difference between the test scores and the training-file scores of
any model can come from the data and not only from the model. For
DeepSeek-R1 7B the two files give almost the same results; for Qwen3.5-27B
they do not. We report the test runs separately
(Sections~\ref{sec:results-qwen} and~\ref{sec:results-deepseek}) and do not
compare them directly with the paper's Table~1.

% DRAFT (to polish): cluster access and queues.
\paragraph{Working on the HPC cluster.}
All our runs need a GPU, so they depend on the university HPC cluster. Two
members of the group had never used this cluster or the Slurm job system
before the project. We had no earlier training in it. We learned it during the
project, from the cluster documentation and with the help of
Inna Emirlova (Special Adviser/HPC administrator at ITU). This
took time that we had planned for the experiments.

The cluster rules make each attempt slow. Package managers and interactive
tools are forbidden on the login node, so every installation and every test
must run as a job. A small error in a job file is found only after the job has
waited in the queue.

The queues were the largest problem. The GPU nodes are shared by many groups.
On 15 November 2026, for example, every GPU of the x86 nodes was in use and
about 40 GPU jobs waited in the queue. Two short test jobs for DeepSeek-R1 7B
did not start that day. So we moved the DeepSeek runs to a less used ARM
node. This solved the waiting time, but it added work: the node has a
different Ollama version, which returns the reasoning text in a different
place, and the Python environment had to be built again for ARM. Someone who
repeats our study should plan for queue time, not only for the GPU hours in
Table~\ref{tab:compute}.

\subsection{Communication with original authors}
We contacted the first author, Dr. Federico Errica, to clarify details of the original experimental setup that were not specified in the paper. In particular, we asked which inference framework and numerical precision or quantization were used for Mixtral-8x7B-Instruct-v0.1 and Llama-3-70B-Instruct. We also reported that the LLM TREC.ipynb notebook in the public repository appeared to be truncated and could not be opened as a valid Jupyter notebook.  

\par\medskip

Dr. Errica responded and consulted the IT personnel responsible for serving the original models. Although the original models were no longer available, they recalled that Llama-3-70B-Instruct was most likely served using INT8 precision and Mixtral-8x7B-Instruct-v0.1 using FP16. However, they could not confirm these settings with certainty. Consequently, we used FP16 for Mixtral, while acknowledging that the exact precision used in the original experiments remains uncertain. 

\par\medskip

Regarding the missing TREC notebook, Dr. Errica acknowledged the issue and recovered an older version of the file, which he subsequently uploaded to the repository. This allowed us to access the previously missing parts of the evaluation code.

\par\medskip

Screenshots documenting the correspondence with the original author
and the repository issue are provided in Appendix~\ref{appendix:communication}.

\subsection{Future work}
\label{sec:future-work}
Testing the explanations in Section~\ref{sec:disc-models} needs more runs
than we could do in this project. Some ideas:

\begin{itemize}
  \item \textbf{Same rephrasings for all models.} Run every model with one
  fixed set of rephrasings as well as its own, as in multi-prompt
  evaluation~\cite{Mizrahi:2024}. This would show how much of the
  difference comes from the rephrasings.
  \item \textbf{More models.} Compare models of the same age but different
  size, and models of the same size from different years.
  \item \textbf{Keep the raw answers.} Save the full model answers, not only
  the extracted class, to see what the N/A answers on RTE say.
  \item \textbf{Older models on the CB test samples.} If Mixtral or Llama
  show the same gap between CB training and test samples, the gap comes from
  the samples and not from contamination. DeepSeek-R1 7B shows no such gap.
  \item \textbf{DeepSeek-R1 7B on all datasets.} Run RTE, DBPedia and WoS,
  ideally with the reasoning switched off if a newer server version allows it.
  \item \textbf{Per-class results on WoS.} Look at the confusion matrices to
  see which topics get mixed up.
  \item \textbf{Repeated runs.} Run some configurations again with other
  precisions or serving settings to see how much the metrics change.
\end{itemize}
